% ==============================================================================
% Model Section
% ==============================================================================


\begin{frame}[t]
    \frametitle{A spatial model with agglomeration and local information}

    \small
    \begin{itemize}

        \only<1->{
        \item \alert<1>{\textbf{Locations}} \(j = 1,\dots,J\)
        \only<1>{
            \begin{itemize}
                \footnotesize
                \item Existing housing stock \(S_j^0\); approved new housing \(q_j\).
                \item Total housing:
                \[
                    S_j = S_j^0 + q_j
                \]
                \item Location fundamentals: amenity \(a_j\), access cost \(\tau_j\), construction cost \(K_j(q_j)\), local development cost \(C_j(q_j)\).
            \end{itemize}
        }
        }

        \only<2->{
        \item \alert<2>{\textbf{Households}} \(g \in \{O,R,M\}\)
        \only<2>{
            \begin{itemize}
                \footnotesize
                \item Owners \(O\), incumbent renters \(R\), potential entrants \(M\).
                \item Entrants are marginal migrants whose location choice is pinned down by an outside option.
                \item Indirect utility:
                \[
                    v_{gj}
                    =
                    a_j
                    +
                    \log w(N)
                    -
                    \alpha_g \log r_j
                    -
                    \tau_j
                    -
                    C_j(q_j)
                    +
                    \varepsilon_{gj}
                \]
                \item \(C_j(q_j)\) captures real local disamenities from development.
                \item In the approval rules below, \(C_j(q_j)\) aggregates these disamenities across affected households.
            \end{itemize}
        }
        }

        \only<3->{
        \item \alert<3>{\textbf{Citywide spillover}}
        \only<3>{
            \begin{itemize}
                \footnotesize
                \item More housing allows more people to live in the city.
                \item Agglomeration:
                \[
                    \log w(N) = \bar w + \gamma \log N
                \]
                \item This creates a diffuse benefit that local districts may not internalize.
            \end{itemize}
        }
        }

        \only<4->{
        \item \alert<4>{\textbf{Homeowners hold local assets}}
        \only<4>{
            \begin{itemize}
                \footnotesize
                \item Owners receive consumption utility and capital gains/losses.
                \[
                    W_j^O = CS_j^O + m_j^O \Delta P_j
                \]
                \item This separates homeowner welfare from renter welfare.
                \item New supply can raise aggregate welfare while reducing local scarcity rents.
            \end{itemize}
        }
        }

        \only<5->{
        \item \alert<5>{\textbf{Land-use authority approves or rejects}} \(q_j\)
        \only<5>{
            \begin{itemize}
                \footnotesize
                \item Developers propose privately feasible projects.
                \item The political regime determines whether projects clear the approval stage.
                \item Key welfare channels: rents, wages, amenities, local costs, and owner asset values.
            \end{itemize}
        }
        }

    \end{itemize}
\end{frame}


% ------------------------------------------------------------
% Slide 2
% ------------------------------------------------------------

\begin{frame}[t]
    \frametitle{Two regimes, three wedges}

    \small
    \begin{itemize}

        \only<1->{
        \item \alert<1>{\textbf{Social benchmark}}
        \only<1>{
            \begin{itemize}
                \footnotesize
                \item A proposed project is socially valuable if total surplus is:
                \[
                    B_j^L + B_j^E - C_j(q_j) - K_j(q_j) \geq 0
                \]
                \item \(B_j^L\): local benefits.
                \item \(B_j^E\): external citywide benefits.
                \item \(C_j(q_j)\): true local cost.
                \item \(K_j(q_j)\): resource/construction cost.
                \item These rules summarize reduced-form project surplus; privately infeasible projects are not proposed.
            \end{itemize}
        }
        }

        \only<2->{
        \item \alert<2>{\textbf{Decentralized approval}}
        \only<2>{
            \begin{itemize}
                \footnotesize
                \item Local politician observes local costs but weighs only local political benefits:
                \[
                    \widetilde{B}_j^L - C_j(q_j) - K_j(q_j) \geq 0
                \]
                \item Strength: local information.
                \item Weakness: ignores citywide and potential-entrant benefits.
            \end{itemize}
        }
        }

        \only<3->{
        \item \alert<3>{\textbf{Centralized approval}}
        \only<3>{
            \begin{itemize}
                \footnotesize
                \item Central authority internalizes citywide benefits but may mismeasure local costs:
                \[
                    B_j^L + B_j^E - \widetilde{C}_j(q_j) - K_j(q_j) \geq 0
                \]
                \item Strength: internalizes spillovers.
                \item Weakness: imperfect local information.
            \end{itemize}
        }
        }

        \only<4->{
        \item \alert<4>{\textbf{Externality wedge:} \(B_j^E\)}
        \only<4>{
            \begin{itemize}
                \footnotesize
                \item Rent reductions outside the district.
                \item Benefits to potential entrants.
                \item Agglomeration/productivity effects.
                \item Fiscal or labor-market benefits outside the approving district.
            \end{itemize}
        }
        }

        \only<5->{
        \item \alert<5>{\textbf{Representation wedge:} \(B_j^L - \widetilde{B}_j^L\)}
        \only<5>{
            \begin{itemize}
                \footnotesize
                \item Local politics may overweight incumbent homeowners.
                \item Renters, potential entrants, and nonlocal beneficiaries may be underweighted.
                \item This is the political-economy channel.
            \end{itemize}
        }
        }

        \only<6->{
        \item \alert<6>{\textbf{Information wedge:} \(C_j(q_j) - \widetilde{C}_j(q_j)\)}
        \only<6>{
            \begin{itemize}
                \footnotesize
                \item Central authority may miss true neighborhood costs.
                \item This is the Jane Jacobs channel \cite{jacobs_life_1961}.
                \item Decentralization can dominate when this wedge is large.
            \end{itemize}
        }
        }

    \end{itemize}

    \only<7>{
        \vspace{0.4em}
        \begin{block}{Regime tradeoff}
            \vspace{.5em}
            \footnotesize
            Centralization solves the externality and representation problems, but may create an information problem. Decentralization uses local information, but may ignore citywide benefits and underrepresented households.
        \end{block}

        \vspace{-0.2em}

    }

\end{frame}


% ------------------------------------------------------------
% Slide 3
% ------------------------------------------------------------

\begin{frame}[t]
    \frametitle{Mapping the model to the empirics}

    \small
    \begin{itemize}

        \only<1->{
        \item \alert<1>{\textbf{Political supply wedge}}
        \only<1>{
            \begin{itemize}
                \footnotesize
                \item Developers build until price equals marginal cost plus the approval wedge:
                \[
                    p_{jt} = K_j'(q_{jt}) + \phi_{jt}
                \]
                \item Model the approval wedge as:
                \[
                    \phi_{jt}
                    =
                    \phi_j^0
                    +
                    \lambda_t \rho H_j
                    +
                    u_{jt}
                \]
                \item \(\rho > 0\): homeowner-heavy districts have higher political costs of approving large projects.
            \end{itemize}
        }
        }

        \only<2->{
        \item \alert<2>{\textbf{Institutional timing}}
        \only<2>{
            \begin{itemize}
                \footnotesize
                \item \(\lambda_t\) captures the strength of local control.
                \item Pre-Morris: final review is less district-specific.
                \item Post-Morris: district-level authority increases.
                \item As member deference hardens:
                \[
                    \lambda_t \uparrow
                \]
                \item Prediction is a growing homeowner gradient, not necessarily an immediate 1990 break.
            \end{itemize}
        }
        }

        \only<3->{
        \item \alert<3>{\textbf{Reduced-form analogue}}
        \only<3>{
            \begin{itemize}
                \footnotesize
                \item Estimate how the homeowner gradient evolves over time:
                \[
                    Y_{j\tau}
                    =
                    \alpha_j
                    +
                    \delta_{b(j),\tau}
                    +
                    \sum_{\ell \neq \tau_0}
                    \beta_\ell
                    \left(
                    H_j \times \mathds{1}\{\tau=\ell\}
                    \right)
                    +
                    X_j'\Gamma_\tau
                    +
                    \varepsilon_{j\tau}
                \]
                \item \(Y_{j\tau}\) is housing production; \(\tau_0\) is the omitted baseline period.
                \item \(b(j)\) is district \(j\)'s assigned majority borough.
                \item \(\beta_\ell\) estimates the period-specific political supply gradient.
            \end{itemize}
        }
        }

        \only<4->{
        \item \alert<4>{\textbf{Predicted margins}}
        \only<4>{
            \begin{itemize}
                \footnotesize
                \item Large buildings should be most affected:
                \[
                    5+ \text{ unit buildings}
                \]
                \item Small buildings should be less affected:
                \[
                    1\text{--}4 \text{ unit buildings}
                \]
                \item The \(1\text{--}4\) unit series is the placebo margin.
            \end{itemize}
        }
        }

        \only<5->{
        \item \alert<5>{\textbf{What the QSM will use}}
        \only<5>{
            \begin{itemize}
                \footnotesize
                \item The reduced form disciplines \(\lambda_t \rho\).
                \item That object becomes the political supply wedge in counterfactuals.
                \item Compare welfare under:
                \begin{enumerate}
                    \scriptsize
                    \item centralized approval,
                    \item decentralized approval,
                    \item decentralized approval with compensation/pricing.
                \end{enumerate}
            \end{itemize}
        }
        }

    \end{itemize}
\end{frame}
